\section{Control de Formaciones en Sistemas No Holonómicos}
\onecolumn
\subsection{Efecto Látigo y Noción de Líder Físico}
Error de formación sobre errores relativos, con atracción a un líder físico:
\begin{equation*}
    \mathbf{e}_i =
    \underbrace{-\sum_{j \in \mathcal{N}_i} a_{ij}\big[(\mathbf{x}_i-\mathbf{r}_i)-(\mathbf{x}_j-\mathbf{r}_j)\big]}_{\text{\footnotesize Consenso entre vecinos}}
    \;\;
    \underbrace{-\,\beta\, b_i^0\big[(\mathbf{x}_i-\mathbf{r}_i)-\mathbf{x}_0\big]}_{\text{\footnotesize Atracción al ancla virtual}}
\end{equation*}

Como cada seguidor persigue una posición referida al líder, un giro brusco de este se propaga en cadena y se amplifica hacia los extremos de la formación: el efecto látigo.
\begin{figure*}[h]
    \centering
    \includegraphics[width=0.9\textwidth]{/home/calafaker/Control-Formation/documentacion/figuras/personales/formacion_v_gazebo.pdf}
    \caption{Formación en V con líder físico en Gazebo: trayectorias y evolución del error de formación de los seguidores en cada giro del líder}
    \label{fig:consensus_results}
\end{figure*}
